% Storyline 1 (final), paragraph 3: the first step is well studied; the second step and the gap.
Most remedies act on the first step. Public institutions, gossip, empathetic evaluation and the synchronization of private opinions reduce how often observers disagree\cite{radzvilavicius2021,kawakatsu2024,radzvilavicius2019,murase2024synchronization,tham2026}, and generous or lenient assessment reduces how often a disputed act is condemned\cite{schmid2021,okada2024}. At the second step, the standard framework lets each new assessment replace an individual's image, so a single disputed act of punishment overturns a standing built over many interactions\cite{ohtsuki2004,hilbe2018}. Quantitative scores, simple records, graded reputations and the aggregation of several observations soften this impact and can make cooperation more robust\cite{schmid2023,clark2020,tanabe2013,murase2022,michelmata2024,yamamoto2025}. Most of them, however, treat reputational change as a single symmetric process, in which positive and negative evidence enter through the same score, threshold or observation window and affect treatment only once a threshold is crossed. Human evaluation departs from both features. Impressions of others are revised asymmetrically, although the direction of the asymmetry depends on context\cite{skowronski1989,slovic1993,mendesiedlecki2013diagnostic,siegel2018}. Institutions, too, grade their response and treat the two directions differently, escalating sanctions only with repeated violations\cite{ostrom1990,cox2010} and restoring an offender's standing only after years without a further offence\cite{blumstein2009}. How the direction of reputational change and the treatment of incomplete evidence affect the punishment dilemma, and with it the evolution of cooperation, remains unclear.
